\begin{abstract}
Agent-based simulations are often too slow for the sweeps a study needs, hard to reproduce,
and hard for a reviewer to check. \modelly{} is an independent Python framework whose
programming model is compatible with Mesa's. It stores agent state in model-owned arrays
(columns) or counted populations (cohorts), runs its inner loops in C with pure-Python twins,
and includes the steps of a study (sweeps, calibration, sensitivity analysis, emulators, data
assimilation, manifests, bundles) and seventeen tested reference models. On eight workloads
written idiomatically for each framework, the column form runs its steps 2.9 to 33 times
faster than per-agent Mesa~3.5 code (12 times faster than Mesa~3.5 with a KD-tree on flocking),
4.6 to 42 times faster than Mesa~2.3 and 3.0 to 9.6 times faster than mesa-frames (on its
four), and 1.8 to 4.5 times slower than hand-written
numpy without agent objects (on three); the cohort form of the one workload that has one runs
135 times faster than Mesa~3.5. Over 100 paired seeds, no difference between a statistic of a
\modelly{} variant and of Mesa~2.3 survives Holm correction, and 63 of 67 comparisons are
equivalent within half a standard deviation (29 with Holm-corrected equivalence tests).
Mesa~3.5 models migrated by a codemod without changing their logic, Mesa's own examples among
them, run 1.05 to 2.3 times faster on \modelly{}, two Mesa~2.3 models 2.7 and 4.8 times. \modelly{} also splits a run
across processes byte-identically, runs columns on a GPU (steps 35 times faster at ten million
agents, 1.5 to 1.9 times including construction), differentiates models, assimilates live
data and exports a model to a web page that runs in any browser.
\end{abstract}
